% !TEX root = ../main.tex
% =====================================================================
\section{Grundlagen}
% =====================================================================
\subsection{Mechanische Grundgrössen}
\begin{symtab}
m & Masse & [kg] \\
s & Strecke & [m] \\
v & Geschwindigkeit & [m/s] \\
a & Beschleunigung & [m/s$^2$] \\
F & Kraft & [N] \\
B & Impuls & [kg*m/s] \\ \midrule
J & Trägheitsmoment & [kg*m$^2$] \\
$\varphi$ & Winkel & [rad] \\
$\omega$ & Winkelgeschwindigkeit & [rad/s] \\
$\omega'$ & Winkelbeschleunigung & [rad/s$^2$] \\
M & Drehmoment & [Nm] \\
D & Drall & [kg*m$^2$/s] \\ \midrule
P & Leistung & [W] \\
E & Energie & [Ws], [Nm], [J] \\
\end{symtab}

Translatorische Bewegung:
\begin{gather*}
v = \frac{ds}{dt} \\
a = \frac{dv}{dt} = \frac{d^2s}{dt^2} \\
F = m \cdot a \\
B = m \cdot v \\
P = F \cdot v \\
E = \int p(t) \cdot dt \\
E_{kin} = \frac{mv^2}{2} \\
E = \int F(s) \cdot ds
\end{gather*}

Rotatorische Bewegung:
\begin{gather*}
\omega = \frac{d\varphi}{dt} = \frac{v}{r} \\
\omega' = \frac{d\omega}{dt} = \frac{d^2\varphi}{dt^2} \\
M = J \cdot \omega' = F \cdot r \\
D = J \cdot \omega \\
E_{kin} = \frac{J\omega^2}{2} \\
E = \int M(\varphi) \cdot d\varphi
\end{gather*}

\columnbreak
% ---------------------------------------------------------------------
\subsection{Magnetismus}
\begin{symtab}
B & Magnetfeld (magnetische Induktion) & [T] \\
H & Feldstärke & [A/m] \\
N & Anz. Wicklungen & [] \\
A & Fläche & [m$^2$] \\
L & Induktivität & [H] \\
$\mu_0$ & Natürliche Permabilität & [Vs/Am] \\
$\mu_r$ & Materialabhängige Permeabilität & [] \\
$\phi$ & Magnetischer Fluss & [Vs] \\
l, l$_{fe}$ & länge, länge Eisenkern & [m] \\
F & Kraft & [N] \\
i, I & Strom & [A] \\
\end{symtab}

\subsubsection{Lorentz-Kraft}
\fig[0.25]{lorentz_kraft}
\begin{gather*}
F = i \cdot l \cdot B
\end{gather*}

\subsubsection{Magnetischer Kreis}
\fig[0.25]{magnetischer_kreis}

Durchflutungsgesetz:
\begin{gather*}
\oint H \cdot ds = H \cdot l_{fe} = \sum I = N \cdot I \\
H = \frac{N \cdot I}{l_{fe}} \\
B = \mu_0 \cdot \mu_r \cdot H \\
\phi = B \cdot A \\
\mu_0 = 4\pi 10^{-7} \\
L = \frac{N^2 \cdot A \cdot \mu_0 \cdot \mu_r}{l_{fe}}
\end{gather*}

\columnbreak
% ---------------------------------------------------------------------
\subsection{C vs. L}
\begin{symtab}
I & Strom & [A] \\
U & Spannung & [V] \\
C & Kapazität & [F] \\
L & Induktivität & [H] \\
$\omega$ & Kreisfrequenz & [rad/s] \\
E & Energie & [Ws] \\
\end{symtab}
\fig[0.9]{c_vs_l}
{\setlength{\jot}{3pt}
\begin{gather*}
E_C = \frac{C \cdot U^2}{2} \\
I_C = C \cdot \frac{du_C}{dt} \\
U_C = \frac{1}{C} \int i_C(t) dt \\
E_L = \frac{L \cdot I^2}{2} \\
I_L = \frac{1}{L} \int u_L(t) dt \\
U_L = L \cdot \frac{di_L}{dt}
\end{gather*}}

% ---------------------------------------------------------------------
\subsection{Leistung}
\begin{symtab}
I & Strom & [A] \\
U & Spannung & [V] \\
P & Leistung & [W] \\
S & Scheinleistung & [VA] \\
Q & Blindleistung & [var] \\
T & Periodendauer & [s] \\
$\lambda$ & Leistungsfaktor & [] \\
$\eta$ & Wirkungsgrad & [\%] \\
cos $\varphi$ & & [] \\
\end{symtab}

\subsubsection{einphasig}
\begin{gather*}
p(t) = u(t) \cdot i(t) \\
P = \frac{1}{T} \int_0^T p(t)\, dt = \frac{1}{T} \int_0^T u(t) \cdot i(t)\, dt \\
S = U_{rms} \cdot I_{rms} \\
P = S \cdot \cos(\varphi) \\
Q = \sqrt{S^2 - P^2} = S \cdot \sin(\varphi) \\
\lambda = \frac{P}{S} \\
\eta = \frac{P_{abgegeben}}{P_{aufgenommen}} = \frac{P_{abgegeben}}{P_{abgegeben} + P_{Verlust}}
\end{gather*}

\subsubsection{dreiphasig}
\begin{gather*}
S = 3 \cdot U_{ph\_rms} \cdot I_{ph\_rms} = \sqrt{3} \cdot U_{vk\_rms} \cdot I_{ph\_rms} \\
\ul{S} = \frac{3}{2} \cdot \ul{U_1} \cdot \ul{I_1^*} \\
P = S \cdot cos(\varphi) \\
Q = \sqrt{S^2 - P^2}
\end{gather*}
